\documentclass[11pt,a4paper]{article}
\usepackage[T1]{fontenc}
\usepackage{newtxtext,newtxmath}
\usepackage{microtype}
\usepackage{xurl}
\usepackage[a4paper,margin=1in]{geometry}
\usepackage{hyperref}

\title{\textbf{AION Phase 22E.47 --- Live Level-3 Post-Game Failure Review}}
\author{Kevin Robinson\\Tessaris AI}
\date{\today}

\begin{document}
\maketitle

\section{Purpose}

Phase 22E.47 records the first live Level-3 test after the Phase 22E.46 dynamic intent switching upgrade.

The test successfully exercised the live stack at Level 3, but AION lost by mate. This document locks the failure map before the next repair phase.

\section{Live Result}

\begin{verbatim}
Game ID: yEGNrLlw
Full ID: yEGNrLlwCg20
Target level: 3
AION colour: white
Final status: mate
Final winner: black
Final ply count: 74
AION move count: 37
Move send OK count: 37
Post-400 recovery count: 0
Trace hash: 0278cdef502c900b5c2a8d005471f97000631628d6f07dd837487f04ceaef15f
\end{verbatim}

\section{Technical Result}

The live execution layer succeeded.

\begin{itemize}
  \item The Level-3 challenge was created.
  \item The network gate was open.
  \item The run was not a dry run.
  \item All AION move posts succeeded.
  \item No post-400 recovery was required.
  \item The 22E.46 live sender stack was active.
\end{itemize}

\section{Observed Failure Map}

\subsection{Wing-Pawn Drift Returned}

After the opening correction, AION again played:

\begin{verbatim}
h2h4
g2g4
\end{verbatim}

These moves were accepted under \texttt{initiative\_pressure}. Against Level 2 this was survivable, but against Level 3 it created long-term king and structure weaknesses.

\subsection{Dynamic Intent Switching Was Active But Unstable}

The new 22E.46 switching layer did activate. It switched between:

\begin{verbatim}
rapid_development
initiative_pressure
center_control
neutralise_critical_threat
king_safety
\end{verbatim}

However, the switching was not yet stable enough. AION often changed labels without enforcing a coherent board plan.

\subsection{Promotion Conversion Failed}

AION successfully created and promoted a passed pawn:

\begin{verbatim}
g6g7
g7g8q
\end{verbatim}

But the promoted queen was then lost:

\begin{verbatim}
g7g8q
...
g1g8
\end{verbatim}

This means queen-first promotion worked, but queen survival and post-promotion conversion were not protected.

\subsection{King Safety Became Passive}

The king-safety mode produced repeated king moves:

\begin{verbatim}
e1f1
f1g1
g1g2
g2f3
f3e2
e2f1
f1e1
e1d1
d1c2
c2b1
b1a1
\end{verbatim}

This avoided immediate threats in places, but it did not remove the opponent's attacking resources or stop the mating net.

\subsection{Endgame Threat Removal Failed}

The final mate occurred after:

\begin{verbatim}
b1a1
c5d4
\end{verbatim}

The Level-3 test exposed that AION needs a king-safety policy that removes the attacking piece or blocks the mating mechanism, not only moves the king.

\section{Next Repair Priorities}

The next repair sequence SHOULD be:

\begin{enumerate}
  \item Phase 22E.48 --- Promotion Safety / Queen Survival Guard.
  \item Phase 22E.49 --- Initiative Hygiene Guard: reject wing-pawn lunges unless concrete.
  \item Phase 22E.50 --- King Safety Threat-Removal Policy.
  \item Phase 22E.51 --- Level-3 Rematch Validation.
\end{enumerate}

\section{Boundary}

This review is evidence-only. It does not call Stockfish, does not use LLM move judgement, and does not use Lichess analysis.

\bigskip

\noindent
Lock ID: AION-PHASE22E47-LIVE-LEVEL3-POST-GAME-FAILURE-REVIEW\\
Status: LOCKED --- Level-3 failure map recorded\\
Maintainer: Tessaris AI\\
Author: Kevin Robinson.

\end{document}
